\documentclass[10pt,a4paper]{amsart}
\usepackage[margin=19mm]{geometry}
\usepackage{amsmath,amssymb,mathtools}
\usepackage[scr=rsfs,cal=euler]{mathalfa}
\usepackage{xfrac}
\usepackage[unicode,hidelinks,bookmarks=false]{hyperref}
\newcommand{\ie}{i.e.\ }
\newcommand{\cf}{cf.\ }
\newcommand{\resp}{respectively\ }
\newcommand{\Subsection}{\subsection}
\providecommand{\nobreakdash}{\nobreak\hbox{-}}
\setlength{\emergencystretch}{2em}
\multlinegap0pt
\usepackage{mathtools}
\usepackage{amssymb}
\usepackage{tikz}
\usepackage{dsfont}
\newtheorem{lemma}{Lemma}
\newcommand{\A}{\mc A}
\newcommand{\mc}{\mathcal}
\newcommand{\sub}{\varrho}
\title{\bf Unique ergodicity, and not, for primitive substitutions on compact alphabets containing isolated points}
\author{James J. Walton}
\date{Source: arXiv:2609.07550v1 (CC BY 4.0)}
\begin{document}
\maketitle

\noindent\emph{Source and licence.} \bf Unique ergodicity, and not, for primitive substitutions on compact alphabets containing isolated points. Version arXiv:2609.07550v1 (CC BY 4.0), distributed by arXiv under the Creative Commons Attribution 4.0 International licence, http://creativecommons.org/licenses/by/4.0/, as recorded in the arXiv licence metadata for this exact version and shown on the abstract page https://arxiv.org/abs/2609.07550v1. Full licence notice: \texttt{licenses/arxiv-2609.07550-CC-BY-4.0.txt}.\par\smallskip

\noindent\textit{Editorial scope.} Counted unit: the article's lemma reduction (source0.tex lines 477--479). The article's complete original proof of it is delivered with it (source0.tex lines 481--511, one \texttt{proof} environment of 31 lines). No mathematics is written, repaired or re-derived: the statement and the proof are the article's own lines, in the article's own notation, and no retained line is substituted or re-flowed.  No further source-local context is needed: every cross-reference of every retained line resolves inside the two delivered spans.  Nothing was omitted from any retained span, so the package carries no omission record.  No retained line contains a citation, so the wrapper carries no bibliography and no bibliography entry.  No retained line contains a figure environment, an image-inclusion command or drawing code, so no image is delivered and the package declares no asset dependency.  Editorial formatting only, in addition: an AMS \texttt{amsart} preamble that restates the article's own declarations and macros used by the retained lines, namely the environments \texttt{lemma} on separate counters, together with the standard packages those lines need.  The retained environments print in this document as Lemma 1 in order of appearance, whereas the article prints the same items with the numbering of its own full document.  The complete native source of the article as distributed in the arXiv e-print for this exact version is bundled unchanged as \texttt{sources/209638/source0.tex}.
\begin{lemma}\label{lem:essential spectral radius unaffected by reduction}
If \(\tau\) is a reduction of \(\sub\), with substitution operators \(M_\tau\) and \(M\), respectively, then \(r_\mathrm{ess}(M_\tau) = r_\mathrm{ess}(M)\).
\end{lemma}

\begin{proof}
We show that \(M_\tau\) is a compact perturbation of \(M\), upon applying any of the three reduction procedures. First, let \(\tau\) be defined by an (\(\varepsilon_\mathrm{in}\)) reduction, for a given finite set \(P\) of isolated points. Let \(Q = \A \setminus P\) be the complementary points and define \(V_P \colon E \to E\) by 
\[
(V_P f)(a) =
\begin{cases}
f(a) & \text{for} \ a \in P \\
0    & \text{otherwise.}
\end{cases}
\]
Note that \(V_P \colon E \to E\), since \(\A = P \sqcup Q\) is a clopen partition. Clearly, \(V_P\) is linear and has image the \(\#P\)-dimensional subspace of functions supported on \(P\), so is finite rank and thus compact. Define \(V \coloneqq V_P \circ M\) which, as a composition of a bounded and compact operator, is also compact. With \(I\) the identity operator, \(I - V_P\) is defined by \((I - V_P)(f) = f |_Q\), where \((f |_Q)(a) = f(a)\) for \(a \in Q\) and \((f |_Q)(a) = 0\) otherwise. Thus,
\[
(M-V)(f) = ((I-V_P) \circ M)(f) = (Mf) |_Q .
\]
It follows that \(M - V = M_\tau\). Indeed, recall that \(\tau\) is defined identically to \(\sub\) on letters in \(Q\) and sends letters in \(P\) to the empty word which, from the operator's perspective, is equivalent to setting \((M_\tau(f))(p) = 0\) for \(p \in P\).

Now suppose that \(\tau\) is instead defined by (\(\varepsilon_\mathrm{out}\)) reduction. The proof is similar to the above: this time we define, instead, \(V \coloneqq M \circ V_P\). Then \((M-V)(f) = (M \circ (I-V_P))(f) = M (f|_Q)\). Thus, \(M-V = M_\tau\), since removing occurrences of letters in \(P\) from superwords has the same effect as setting \(f\) equal to \(0\) on such letters and then applying the substitution operator.

Finally, let \(\tau\) be defined by (\(\varepsilon_\mathrm{con}\)) reduction, `erasing a constant', with respect to a clopen set \(U\) on which each superword contains an occurrence of \(b \in \A\). Let \(\chi_U\) be the indicator function of \(U\), that is \(\chi_U(a) = 1\) for \(a \in U\) and \(\chi_U(a) = 0\) for \(a \notin U\). Since \(U\) is clopen, we have \(\chi_U \in E\). Then we define the operator \(V_U^b \colon E \to E\) as
\[
V_U^b(f) \coloneqq (f(b)) \cdot \chi_U .
\]
Clearly, \(V_U^b\) is linear and has rank one, since its image is spanned by \(\chi_U \in E\) (so, in particular, \(V_U^b\) is compact). Then \((V_U^b(f))(a) = 0\) if \(a \notin U\) and \((V_U^b(f))(a) = f(b)\) otherwise, thus
\[
((M - V_U^b)f)(a) =
\begin{cases}
(Mf)(a)        & \text{for } a \notin U \\
(Mf)(a) - f(b) & \text{for } a \in U . \\
\end{cases}
\]
It follows that \(M-V_U^b = M_\tau\), since the above has had the effect of removing one occurrence of \(b\) from the substitution of any letter in \(U\). Since applying any finite combination of the three reduction procedures defines a new operator that only differs by the addition of a compact operator, the essential spectral radius remains invariant.
\end{proof}

\end{document}
